% ---- pg-579 （书页 567） ----

\noindent\textbf{例 12.19}\quad (大学及其特征的双重信息图)

\hyperref[tab:12.9]{表 12.9} 给出某些大学对一些变量的数据, 用于将主要大学进行比较或排队. 这些变量包括 $X_1$=新生的平均 SAT 得分, $X_2$=新生在高中时期名列班上 前 $10\%$ 的人数百分比, $X_3$=报名者被接受入学的百分比, $X_4$=学生与教师的比例, $X_5$=估计的年费用, 及 $X_6$=毕业比例($\%$).

\begin{mvtable}
\phantomsection\label{tab:12.9}\mvtabcap{表 12.9\quad 大学数据}
\setlength{\tabcolsep}{3.5pt}%
\zihao{6}%
\begin{longtable}{lrrrrrr}\toprule
\textbf{大\quad 学} & $X_1$ & $X_2$ & $X_3$ & $X_4$ & $X_5$ & $X_6$ \\
\midrule
\endfirsthead
\multicolumn{7}{c}{\zihao{6}续表}\\
\toprule
\textbf{大\quad 学} & $X_1$ & $X_2$ & $X_3$ & $X_4$ & $X_5$ & $X_6$ \\
\midrule
\endhead
\bottomrule
\endlastfoot

哈佛大学(Harvard) & 14.00 & 91 & 14 & 11 & 39{,}525 & 97 \\
普林斯顿大学(Princeton) & 13.75 & 91 & 14 & 8 & 30{,}220 & 95 \\
耶鲁大学(Yale) & 13.75 & 95 & 19 & 11 & 43{,}514 & 96 \\
斯坦福大学(Stanford) & 13.60 & 90 & 20 & 12 & 36{,}450 & 93 \\
麻省理工学院(MIT) & 13.80 & 94 & 30 & 10 & 34{,}870 & 91 \\
杜克大学(Duke) & 13.15 & 90 & 30 & 12 & 31{,}585 & 95 \\
加州理工学院(CalTech) & 14.15 & 100 & 25 & 6 & 63{,}575 & 81 \\
达特茅斯大学(Dartmouth) & 13.40 & 89 & 23 & 10 & 32{,}162 & 95 \\
布朗大学(Brown) & 13.10 & 89 & 22 & 13 & 22{,}704 & 94 \\
琼斯霍普金斯大学(JohnsHopkins) & 13.05 & 75 & 44 & 7 & 58{,}691 & 87 \\
芝加哥大学(UChicago) & 12.90 & 75 & 50 & 13 & 38{,}380 & 87 \\
宾夕法尼亚大学(UPenn) & 12.85 & 80 & 36 & 11 & 27{,}553 & 90 \\
康奈尔大学(Cornell) & 12.80 & 83 & 33 & 13 & 21{,}864 & 90 \\
西北大学(Northwestern) & 12.60 & 85 & 39 & 11 & 28{,}052 & 89 \\
哥伦比亚大学(Columbia) & 13.10 & 76 & 24 & 12 & 31{,}510 & 88 \\
诺特丹大学(NotreDame) & 12.55 & 81 & 42 & 13 & 15{,}122 & 94 \\
弗吉尼亚大学(UVirginia) & 12.25 & 77 & 44 & 14 & 13{,}349 & 92 \\
乔治敦大学(Georgetown) & 12.55 & 74 & 24 & 12 & 20{,}126 & 92 \\
卡耐基梅隆大学(CarnegieMellon) & 12.60 & 62 & 59 & 9 & 25{,}026 & 72 \\
密歇根大学(UMichigan) & 11.80 & 65 & 68 & 16 & 15{,}470 & 85 \\
加州大学(伯克利分校)(UCBerkeley) & 12.40 & 95 & 40 & 17 & 15{,}140 & 78 \\
威斯康星大学(UWisconsin) & 10.85 & 40 & 69 & 15 & 11{,}857 & 71 \\
宾州大学(PennState) & 10.81 & 38 & 54 & 18 & 10{,}185 & 80 \\
普度大学(Purdue) & 10.05 & 28 & 90 & 19 & 9{,}066 & 69 \\
得克萨斯 A\&M 大学(TexasA\&M) & 10.75 & 49 & 67 & 25 & 8{,}704 & 67 \\
\end{longtable}
\end{mvtable}

{\zihao{6}资料来源: 美国新闻与世界报道, 1995 年 9 月 18 日.}

由于其中两个变量(SAT 得分和费用)在尺度上与其他变量有很大差别, 我们将数据标准化, 并根据标准化观测值 $\mathbf{z}_j$ 的矩阵作双重信息图. 这个双重信息图在\hyperref[fig:12.24]{图 12.24} 中给出.

注意加州理工学院和琼斯霍普金斯大学远离群体单独在一起, 费用变量($X_5$)对这种定位负主要责任. 在我们的样本中, 大型州立大学在图的左侧, 而大部分私立大学集中在图的右侧. 变量 SAT($X_1$)、前 10 名比例($X_2$)和毕业率($X_6$)的较大值是与私立学校组相联系的. 西北大学位于图的中心.\hfill$\blacksquare$

高尔和韩德(Gower \& Hand, 见文献[12])开发了一种新的双重信息图方法, 有某些优点. 他们的方法是对散点图的推广, 具有易于解释的特点:
